\documentclass[12pt,a4paper]{article}

% --- Encoding ---
\usepackage[utf8]{inputenc}
\usepackage[T1]{fontenc}
\usepackage[english]{babel}

% --- Math ---
\usepackage{amsmath, amssymb, mathtools}

% --- Layout ---
\usepackage{geometry}
\geometry{margin=2.5cm, top=3cm, bottom=3cm}

% --- Figures ---
\usepackage{graphicx}
\usepackage{float}
% Figures are generated by notebook_implementation.ipynb and saved to
% METAHEURISTICS/tabu_search/figures/ — compile from docs/src/ or adjust the path below.
\graphicspath{{../../../METAHEURISTICS/tabu_search/figures/}}

% --- Tables ---
\usepackage{booktabs}
\usepackage{tabularx}
\usepackage{array}
\usepackage{multirow}
\usepackage[table]{xcolor}

% --- Code listings ---
\usepackage{listings}
\definecolor{codegray}{rgb}{0.45,0.45,0.45}
\definecolor{codeblue}{rgb}{0.0,0.25,0.65}
\definecolor{codeorange}{rgb}{0.75,0.35,0.0}
\definecolor{codered}{rgb}{0.6,0.0,0.0}
\lstdefinestyle{pythonstyle}{
  language=Python,
  basicstyle=\ttfamily\small,
  keywordstyle=\color{codeblue}\bfseries,
  commentstyle=\color{codegray}\itshape,
  stringstyle=\color{codeorange},
  numberstyle=\tiny\color{codegray},
  showstringspaces=false,
  breaklines=true,
  frame=single,
  rulecolor=\color{gray!40},
  backgroundcolor=\color{gray!5},
  xleftmargin=1em,
  numbers=left,
  numbersep=6pt,
}
\lstdefinestyle{bashstyle}{
  language=bash,
  basicstyle=\ttfamily\small,
  keywordstyle=\color{codered},
  commentstyle=\color{codegray}\itshape,
  showstringspaces=false,
  breaklines=true,
  frame=single,
  rulecolor=\color{gray!40},
  backgroundcolor=\color{gray!5},
  xleftmargin=1em,
}

% --- Hyperlinks ---
\usepackage{hyperref}
\hypersetup{colorlinks=true, linkcolor=blue!70!black, citecolor=blue, urlcolor=blue!80!black}

% --- Headers ---
\usepackage{fancyhdr}
\pagestyle{fancy}
\fancyhf{}
\rhead{TSP --- Metaheuristics}
\lhead{Tabu Search (TS)}
\cfoot{\thepage}

% --- Misc ---
\usepackage{enumitem}

% =============================================================================
% TITLE
% =============================================================================
\title{%
  \LARGE\textbf{Tabu Search (TS)}\\[0.6em]
  \large Applied to the Travelling Salesman Problem (TSP)
}
\date{\today}

% =============================================================================
\begin{document}
% =============================================================================

\maketitle
\thispagestyle{empty}

% --- Revision table ---
\begin{center}
\rowcolors{2}{gray!8}{white}
\begin{tabularx}{\linewidth}{ccX}
\toprule
\textbf{Version} & \textbf{Date} & \textbf{Description} \\
\midrule
1.0 & 2026-01-20 & Initial formalisation: problem setting, algorithm, complexity, benchmarks \\
2.0 & 2026-04-10 & TSPTW-D extensions, full-width tables, symbol annotations, benchmark restructure \\
\bottomrule
\end{tabularx}
\end{center}

\begin{abstract}
This document presents the application of Tabu Search (TS) to the
Travelling Salesman Problem (TSP).
After a formal definition and a complete mathematical description of the
algorithm (Glover, 1986--1989), we describe the search procedure, complexity
analysis, strengths and limitations of the approach.
The benchmark covers sizes $n \in \{10, 50, 100\}$ and reports solution
quality (average tour length, gap relative to the nearest-neighbour baseline),
and average runtime.
The scalability ceiling --- where the $O(n^2)$ neighbourhood evaluation becomes
the bottleneck --- is analysed and quantified.
\end{abstract}

\tableofcontents

\section*{Acronyms}
\addcontentsline{toc}{section}{Acronyms}

\begin{center}
\rowcolors{2}{gray!8}{white}
\begin{tabularx}{\linewidth}{lX}
\toprule
\textbf{Acronym} & \textbf{Definition} \\
\midrule
TS      & Tabu Search \\
TSP     & Travelling Salesman Problem \\
TSPTW   & TSP with Time Windows \\
TSPTW-D & TSP with Time Windows and Disruptions \\
NN      & Nearest Neighbour \\
GNN     & Graph Neural Network \\
FIFO    & First In, First Out \\
LKH     & Lin--Kernighan--Helsgott \\
\bottomrule
\end{tabularx}
\end{center}

\newpage

% =============================================================================
\section{Definition}
% =============================================================================

\textbf{Tabu Search (TS)} is a metaheuristic local search algorithm that
navigates the solution space by iteratively applying moves to a current
solution and using an explicit \emph{memory structure} --- the \textbf{tabu
list} --- to forbid recently visited moves, thereby escaping local optima and
encouraging exploration.

\subsection*{Origin}

\begin{center}
\rowcolors{2}{gray!8}{white}
\begin{tabularx}{\linewidth}{cX}
\toprule
\textbf{Year} & \textbf{Contribution} \\
\midrule
1986 & Glover --- first formulation of TS principles
       (\textit{Computers \& Operations Research}) \\
1989 & Glover --- Tabu Search Part~I, full algorithmic framework
       (\textit{ORSA Journal on Computing}) \\
1990 & Glover --- Tabu Search Part~II, intensification and diversification \\
1997 & Glover \& Laguna --- \textit{Tabu Search} monograph (standard reference) \\
\bottomrule
\end{tabularx}
\end{center}

The name \textit{tabu} derives from the Polynesian word for ``forbidden'',
reflecting the central mechanism: recent moves are declared off-limits to
prevent the search from cycling.

\subsection*{Category}

Tabu Search belongs to the family of \textbf{metaheuristics} --- high-level,
problem-independent strategies that guide an underlying local search to explore
the solution space more effectively than naive descent.
More specifically, TS is a \textbf{local search with memory}: it extends
classical neighbourhood-based local search by maintaining an explicit record
of recent moves.

\subsection*{Key concepts}

\begin{itemize}[leftmargin=1.5em]
  \item \textbf{Tabu list} $T$: a fixed-length FIFO queue of recently performed
    moves, declared \emph{forbidden} for the next $\theta$ iterations (tabu
    tenure).
  \item \textbf{Aspiration criterion}: overrides the tabu status of a move when
    accepting it yields a solution better than the best known so far.
  \item \textbf{Non-improving moves}: unlike steepest descent, TS accepts the
    best available move even if it worsens the objective, allowing escape from
    local optima.
  \item \textbf{Short-term memory}: the tabu list itself --- prevents immediate
    reversal and cycling over the last $\theta$ moves.
  \item \textbf{Long-term memory} (optional): frequency- or recency-based
    structures for \emph{intensification} (exploiting promising regions) and
    \emph{diversification} (forcing exploration of unvisited regions).
\end{itemize}

% =============================================================================
\section{Formal Description}
% =============================================================================

\subsection{Problem setting}

Let $G = (V, E)$ be a complete undirected graph where
$V = \{1, 2, \ldots, n\}$ is the set of cities and $d_{ij} \geq 0$ is the
Euclidean distance between cities $i$ and $j$.
The TSP seeks a Hamiltonian tour (a permutation $\pi$ of $V$) that minimises
the total length:

\begin{equation}
f(\pi) = \sum_{k=1}^{n} d_{\pi(k),\; \pi(k \bmod n\, +\, 1)}
\label{eq:tsp_obj}
\end{equation}

{\small\begin{tabularx}{\linewidth}{@{\hspace{1.5em}}l@{\;---\;}X}
$f(\pi)$ & total Euclidean length of tour $\pi$ \\
$\pi(k)$ & city index visited at step $k$ \\
$d_{ij}$ & Euclidean distance between cities $i$ and $j$ \\
$n$ & number of cities \\
\end{tabularx}}\smallskip

The solution space is $\mathcal{S} = \{\pi \in \mathfrak{S}_n\}$, with
$|\mathcal{S}| = (n-1)!/2$ for the symmetric TSP.

\subsection{Neighbourhood structure --- 2-opt}

Given a current tour $s \in \mathcal{S}$, the \textbf{2-opt neighbourhood}
$N(s)$ contains all tours obtainable by removing two non-adjacent edges
$(s(i), s(i{+}1))$ and $(s(j), s(j{+}1))$ and reconnecting the four endpoints
in the only alternative way (reversing the segment between $i{+}1$ and $j$):

\begin{equation}
N(s) = \bigl\{\, s' \in \mathcal{S} \;\big|\;
  s' = \mathrm{2\text{-}opt}(s, i, j),\;
  1 \leq i < j \leq n,\; j \neq i{+}1 \,\bigr\}
\label{eq:neighbourhood}
\end{equation}

{\small\begin{tabularx}{\linewidth}{@{\hspace{1.5em}}l@{\;---\;}X}
$N(s)$ & 2-opt neighbourhood of tour $s$: all tours reachable by one edge swap \\
$s'$ & candidate tour obtained by reversing the segment between positions $i+1$ and $j$ \\
$\mathcal{S}$ & solution space (all Hamiltonian permutations of $V$) \\
$i, j$ & swap indices; $j \neq i+1$ excludes adjacent (trivial) pairs \\
\end{tabularx}}\smallskip

The \textbf{cost delta} of a 2-opt move $(i, j)$ is computed in $O(1)$ using
the pre-computed distance matrix:

\begin{equation}
\Delta f(s, i, j) =
  \bigl(d_{s(i),\,s(j)} + d_{s(i+1),\,s(j+1)}\bigr)
  - \bigl(d_{s(i),\,s(i+1)} + d_{s(j),\,s(j+1)}\bigr)
\label{eq:delta_2opt}
\end{equation}

{\small\begin{tabularx}{\linewidth}{@{\hspace{1.5em}}l@{\;---\;}X}
$\Delta f(s, i, j)$ & signed cost change when applying 2-opt move $(i,j)$; negative implies improvement \\
$d_{s(i),\,s(j)}$ & distance between the cities at tour positions $i$ and $j$ (new edges added) \\
$d_{s(i),\,s(i+1)}$ & length of the removed edge between positions $i$ and $i+1$ \\
$d_{s(j),\,s(j+1)}$ & length of the removed edge between positions $j$ and $j+1$ \\
\end{tabularx}}\smallskip

A move is \emph{improving} when $\Delta f < 0$.

\subsection{Tabu list}

The tabu list $T$ is a FIFO queue of bounded capacity $\theta$ (the
\emph{tabu tenure}) that stores recently performed moves.
A move $m = (i, j)$ is \textbf{tabu} if $m \in T$:

\begin{equation}
T \subset \mathcal{M}, \qquad |T| \leq \theta
\label{eq:tabu_list}
\end{equation}

{\small\begin{tabularx}{\linewidth}{@{\hspace{1.5em}}l@{\;---\;}X}
$T$ & tabu list: FIFO queue of recently performed moves, declared forbidden \\
$\mathcal{M}$ & set of all possible moves \\
$\theta$ & tabu tenure: number of iterations a move remains forbidden \\
\end{tabularx}}\smallskip

where $\mathcal{M}$ is the set of all possible moves.
A hash set mirrors $T$ to provide $O(1)$ membership tests.

\subsection{Aspiration criterion}

The aspiration criterion $A(s')$ overrides the tabu status of a move when
accepting it leads to a \textbf{global improvement}:

\begin{equation}
A(s') = \mathbf{1}\!\left[\, f(s') < f(s^*) \,\right]
\label{eq:aspiration}
\end{equation}

{\small\begin{tabularx}{\linewidth}{@{\hspace{1.5em}}l@{\;---\;}X}
$A(s')$ & aspiration criterion: Boolean override of tabu status \\
$f(s')$ & objective value (tour length) of candidate solution $s'$ \\
$s^*$ & best solution (shortest tour) found so far \\
$\mathbf{1}[\cdot]$ & indicator function: equals 1 if the condition holds, 0 otherwise \\
\end{tabularx}}\smallskip

where $s^*$ is the best solution found so far.

\subsection{Move selection rule}

At each iteration, the algorithm selects the best admissible move:

\begin{equation}
s^{\mathrm{next}} = \arg\min_{s' \in N(s)}
  \bigl\{\, f(s') \;\big|\;
    \mathrm{move}(s \to s') \notin T \;\lor\; A(s') \,\bigr\}
\label{eq:move_selection}
\end{equation}

{\small\begin{tabularx}{\linewidth}{@{\hspace{1.5em}}l@{\;---\;}X}
$s^{\mathrm{next}}$ & solution selected for the next iteration \\
$N(s)$ & 2-opt neighbourhood of the current tour $s$ \\
$\mathrm{move}(s \to s') \notin T$ & move is not forbidden (not in tabu list) \\
$A(s')$ & aspiration criterion is satisfied (global improvement over $s^*$) \\
\end{tabularx}}\smallskip

\subsection{Tabu list update}

After selecting move $m^* = (i, j)$:
\begin{enumerate}[leftmargin=2em]
  \item Append $m^*$ to $T$.
  \item If $|T| > \theta$, evict the oldest entry (FIFO discipline).
  \item Update $s^*$ if $f(s^{\mathrm{next}}) < f(s^*)$.
\end{enumerate}

This ensures $m^*$ is forbidden for the next $\theta$ iterations,
preventing immediate reversal.

\subsection{Complexity}

\begin{center}
\rowcolors{2}{gray!8}{white}
\begin{tabularx}{\linewidth}{Xccc}
\toprule
\textbf{Component} & \textbf{Time} & \textbf{Space} & \textbf{Dominant term} \\
\midrule
Distance matrix (pre-computed once) & $O(n^2)$ & $O(n^2)$ & --- \\
Neighbourhood evaluation (per iter.) & $O(n^2)$ & $O(1)$ & 2-opt pairs \\
Tabu check (hash set, per move) & $O(1)$ & $O(\theta)$ & --- \\
Apply 2-opt move & $O(n)$ & $O(n)$ & segment reversal \\
\midrule
\textbf{Total} & $\boldsymbol{O(\mathrm{iter} \times n^2)}$ &
  $\boldsymbol{O(n^2 + \theta)}$ & \textbf{neighbourhood} \\
\bottomrule
\end{tabularx}
\end{center}

The number of candidate moves per iteration is exactly
$\binom{n}{2} - (n-1) = \tfrac{(n-1)(n-2)}{2} = O(n^2)$.
Early stopping via a patience counter typically reduces the effective
iteration count well below \texttt{max\_iter}.

% =============================================================================
\section{Architecture / Algorithm Steps}
% =============================================================================

\begin{center}
\begin{tabularx}{\linewidth}{X}
\textbf{Input} city coordinates $(x_1, \dots, x_n)$\\
$\quad\downarrow$\\
\textbf{Step 1} --- Pre-compute distance matrix $D \in \mathbb{R}^{n \times n}$\\
$\quad\downarrow$\\
\textbf{Step 2} --- Initialise tour $s$ via nearest-neighbour heuristic\\
$\quad\downarrow$\\
\textbf{Step 3} --- Repeat until stopping criterion:\\
$\quad\quad\quad$ (a) Evaluate all 2-opt neighbours; compute $\Delta f(s, i, j)$ in $O(1)$\\
$\quad\quad\quad$ (b) Select best admissible move (not tabu, or passes aspiration)\\
$\quad\quad\quad$ (c) Apply move; update tabu list $T$ (FIFO, capacity $\theta$)\\
$\quad\quad\quad$ (d) Update $s^*$ if improved; increment patience counter otherwise\\
$\quad\downarrow$\\
\textbf{Output} best tour $s^*$ and its length $f(s^*)$\\
\end{tabularx}
\end{center}

\textbf{Stopping criteria:}
\begin{itemize}[leftmargin=1.5em]
  \item Hard limit: \texttt{max\_iter} iterations reached.
  \item Early stopping: no improvement for \texttt{patience} consecutive
    iterations.
\end{itemize}

\textbf{Initialisation:}
The nearest-neighbour (NN) heuristic builds an initial tour greedily ---
starting from city~0, repeatedly visiting the closest unvisited city.
This produces a tour typically $20$--$25\,\%$ above optimum on random
instances, giving TS a reasonable warm start.

% =============================================================================
\section{Strengths and Limitations}
% =============================================================================

\subsection*{Strengths}

\begin{center}
\rowcolors{2}{gray!8}{white}
\begin{tabularx}{\linewidth}{lX}
\toprule
\textbf{Property} & \textbf{Detail} \\
\midrule
Local-optima escape & Accepts non-improving moves, unlike steepest descent or
  plain 2-opt \\
No training data required & Can be applied to any TSP instance immediately,
  without labelled solutions \\
Reproducibility & Fully deterministic given a fixed seed and initial solution \\
Simplicity & Conceptually straightforward; few critical hyperparameters
  ($\theta$, \texttt{max\_iter}, \texttt{patience}) \\
Flexibility & Easily extended with long-term memory, candidate lists, or
  problem-specific move operators \\
\bottomrule
\end{tabularx}
\end{center}

\subsection*{Limitations}

\begin{center}
\rowcolors{2}{gray!8}{white}
\begin{tabularx}{\linewidth}{lX}
\toprule
\textbf{Limitation} & \textbf{Impact} \\
\midrule
Quadratic neighbourhood & $O(n^2)$ per iteration; impractical beyond
  $n \approx 1\,000$ without candidate lists \\
Tenure sensitivity & $\theta$ too small $\Rightarrow$ cycling;
  $\theta$ too large $\Rightarrow$ over-restricted search \\
No optimality certificate & Heuristic with no proof of optimality;
  quality depends on landscape and parameters \\
Slow inference vs.\ DL & Seconds per instance at $n = 100$; orders of
  magnitude slower than a trained GNN at inference \\
No quality self-report & Always outputs a tour with no internal confidence
  estimate --- silent failures possible \\
\bottomrule
\end{tabularx}
\end{center}

% =============================================================================
\section{Use Case Explanation}
% =============================================================================

\subsection*{When to use Tabu Search for TSP}

\begin{enumerate}[leftmargin=2em]
  \item \textbf{Medium-scale instances ($n = 50$--$500$).}
    TS reliably closes a large fraction of the NN gap (typically reaching
    5--15\,\% above optimum) in a practical time budget.
    This is its natural sweet spot.

  \item \textbf{When significantly better than NN is needed without labelled
    data.}
    Unlike GNN-based solvers, TS requires no training dataset of solved
    instances.
    It can be deployed immediately on any instance distribution.

  \item \textbf{As a strong, reproducible baseline.}
    In research comparing classical and learned solvers, TS provides a
    deterministic, well-understood reference point that does not depend on
    training quality or random seeds at inference.

  \item \textbf{When moderate computational budget is acceptable.}
    A TS run on $n = 100$ with 5\,000 iterations (patience~500) completes
    in under a second on a standard CPU.
\end{enumerate}

\subsection*{When NOT to use Tabu Search}

\begin{itemize}[leftmargin=1.5em]
  \item \textbf{Very large instances ($n > 5\,000$):} The $O(n^2)$ per-iteration
    cost becomes prohibitive.
    Candidate-list TS, LKH-3, or population-based methods are more appropriate.
  \item \textbf{When a proof of optimality is required:} Use exact solvers
    (Concorde, branch-and-cut).
  \item \textbf{When batch inference speed dominates:} A trained GNN or
    Attention Model can solve thousands of instances per second on a GPU;
    TS cannot match this throughput.
\end{itemize}

\subsection*{Comparison with other TSP solvers}

\begin{center}
\rowcolors{2}{gray!8}{white}
\begin{tabularx}{\linewidth}{lccX}
\toprule
\textbf{Method} & \textbf{Gap at $n=100$} & \textbf{Inference time} &
  \textbf{Training data?} \\
\midrule
Nearest Neighbour                  & ${\sim}20$--$25\,\%$ & ${<}1$\,ms  & No  \\
\textbf{Tabu Search (2-opt)}       & ${\sim}5$--$12\,\%$  & $0.1$--$1$\,s & No \\
Pointer Network (Vinyals 2015)     & ${\sim}8$--$10\,\%$  & ${<}10$\,ms  & Yes \\
Attention Model greedy (Kool 2019) & ${\sim}3$--$5\,\%$   & ${<}10$\,ms  & Yes \\
Attention Model beam search        & ${\sim}1$--$2\,\%$   & ${\sim}100$\,ms & Yes \\
LKH-3 (Helsgott)                   & ${<}0.5\,\%$         & $1$--$10$\,s & No  \\
Concorde (exact)                   & $0\,\%$              & $1$--$60$\,s & No  \\
\bottomrule
\end{tabularx}
\end{center}

TS occupies a middle ground: better than pure NN, competitive with early
learned solvers at greedy decoding, yet no training data is required.

% =============================================================================
\section{Code Implementation}
% =============================================================================

The full solver is defined in \texttt{solver.py}.
The snippet below illustrates solver creation, a forward run, and inspection
of the convergence history.

\begin{lstlisting}[style=pythonstyle,
  caption={Creation, run, and result inspection
           (\texttt{solver.py} + \texttt{data.py})}]
import torch
from solver import TabuSearch
from data   import random_instance, tour_length, load_cities

# --- Random instance (or load from dataset) --------------------------------
n      = 50
coords = random_instance(n, seed=42)   # torch.Tensor (n, 2) in [0, 1]
# coords = load_cities(n, source="tsp")  # from TSP chunk files

# --- Create solver ---------------------------------------------------------
ts = TabuSearch(
    coords,
    tabu_tenure = 20,    # moves stay forbidden for 20 iterations
    max_iter    = 5000,  # hard iteration limit
    patience    = 500,   # stop early if no improvement for 500 iterations
    seed        = 42,
)

# --- Run -------------------------------------------------------------------
best_tour, best_length, history = ts.run()

print(f"Tour length : {best_length:.4f}")
print(f"Iterations  : {len(history)}")
print(f"Tour        : {best_tour}")

# history is a list of best lengths per iteration -- useful for plotting
# convergence: plt.plot(history)
\end{lstlisting}

\begin{lstlisting}[style=bashstyle,
  caption={Running the benchmark notebook from the command line}]
# Launch Jupyter and open the implementation notebook
cd METAHEURISTICS/tabu_search/
jupyter notebook notebook_implementation.ipynb

# Or execute it non-interactively (requires nbconvert)
jupyter nbconvert --to notebook --execute notebook_implementation.ipynb \
        --output notebook_implementation_executed.ipynb
\end{lstlisting}

Key parameters and their recommended ranges:

\begin{center}
\rowcolors{2}{gray!8}{white}
\begin{tabularx}{\linewidth}{lccX}
\toprule
\textbf{Parameter} & \textbf{Default} & \textbf{Typical range} &
  \textbf{Effect} \\
\midrule
\texttt{tabu\_tenure} $\theta$ & 20 & $\lfloor\sqrt{n}\rfloor$--$n/5$ &
  Controls anti-cycling strength; see Section~\ref{sec:experimental} \\
\texttt{max\_iter}   & 5\,000 & 1\,000--50\,000 &
  Hard iteration cap \\
\texttt{patience}    & 500 & 100--2\,000 &
  Early stopping; set to \texttt{max\_iter} to disable \\
\bottomrule
\end{tabularx}
\end{center}

% =============================================================================
\section{TSP Benchmark}
% =============================================================================

The benchmark was run with the medium configuration
(\texttt{tabu\_tenure}$=20$, \texttt{max\_iter}$=5\,000$,
\texttt{patience}$=500$, \texttt{seed}$=42$).
Sizes $n \in \{10, 50, 100\}$ were evaluated on instances loaded from the TSP
chunk dataset.
Results are populated by \texttt{benchmark\_tabu\_tsp.ipynb}.

\subsection*{Metrics}

\begin{center}
\rowcolors{2}{gray!8}{white}
\begin{tabularx}{\linewidth}{lX}
\toprule
\textbf{Metric} & \textbf{Definition} \\
\midrule
Avg.\ length & Mean tour length over all test instances \\
Gap (\%) & $\frac{f_{\mathrm{TS}} - f_{\mathrm{NN}}}{f_{\mathrm{NN}}} \times 100$;
  negative means TS \emph{beats} the NN baseline \\
Avg.\ time (s) & Mean wall-clock time per TS run (nearest-neighbour init
  + main loop) \\
\bottomrule
\end{tabularx}
\end{center}

\textbf{Reference tour:}
nearest-neighbour heuristic (itself ${\sim}20\,\%$ above optimal on random
Euclidean instances); a negative gap indicates TS outperforms this baseline.
For $n \leq 10$, exact optimal tours are additionally reported.

\subsection*{Results --- $n = 10$}

\begin{center}
\rowcolors{2}{gray!8}{white}
\begin{tabularx}{\linewidth}{Xrrr}
\toprule
\textbf{Method} & \textbf{Avg.\ length} & \textbf{Gap (\%)} & \textbf{Avg.\ time (s)} \\
\midrule
Nearest Neighbour     & 3.0469 & 0.0\,(ref.)  & 0.0001 \\
\textbf{Tabu Search}  & 2.7703 & $-9.1$      & 0.071  \\
Optimal (exact)       & 2.7703 & $-9.1$      & ---    \\
\bottomrule
\end{tabularx}
\end{center}

\subsection*{Results --- $n = 50$}

\begin{center}
\rowcolors{2}{gray!8}{white}
\begin{tabularx}{\linewidth}{Xrrr}
\toprule
\textbf{Method} & \textbf{Avg.\ length} & \textbf{Gap (\%)} & \textbf{Avg.\ time (s)} \\
\midrule
Nearest Neighbour    & 7.1567 & 0.0\,(ref.) & 0.0004 \\
\textbf{Tabu Search} & 5.8282 & $-18.6$     & 5.791  \\
\bottomrule
\end{tabularx}
\end{center}

\subsection*{Results --- $n = 100$}

\begin{center}
\rowcolors{2}{gray!8}{white}
\begin{tabularx}{\linewidth}{Xrrr}
\toprule
\textbf{Method} & \textbf{Avg.\ length} & \textbf{Gap (\%)} & \textbf{Avg.\ time (s)} \\
\midrule
Nearest Neighbour    & 9.8126 & 0.0\,(ref.) & 0.0006 \\
\textbf{Tabu Search} & 8.0161 & $-18.3$     & 30.618 \\
\bottomrule
\end{tabularx}
\end{center}

\subsection*{Figures}

\begin{figure}[H]
  \centering
  \includegraphics[width=\linewidth]{benchmark_summary.png}
  \caption{
    Left: gap relative to the nearest-neighbour baseline per instance size
    (negative = TS beats NN).
    Centre: total runtime as a function of $n$.
    Right: convergence curves normalised to the NN initial length for all
    benchmark sizes.
  }
  \label{fig:benchmark_summary}
\end{figure}

\begin{figure}[H]
  \centering
  \includegraphics[width=\linewidth]{benchmark_tours.png}
  \caption{
    Tour visualisations for $n = 10$ and $n = 50$:
    Tabu Search solution (blue, left column) vs.\ nearest-neighbour
    baseline (red, right column).
    The gold dot marks the start city.
  }
  \label{fig:benchmark_tours}
\end{figure}

\begin{figure}[H]
  \centering
  \includegraphics[width=\linewidth]{demo_overview.png}
  \caption{
    Overview on an $n = 10$ instance:
    (i)~convergence curve of the TS best-tour length,
    (ii)~brute-force optimal tour (green),
    (iii)~Tabu Search tour (blue) with gap to optimal.
  }
  \label{fig:demo_overview}
\end{figure}

% =============================================================================
\subsection{Scalability analysis for large instances}
\label{sec:scalability}
% =============================================================================

The dominant cost per iteration is the full 2-opt neighbourhood evaluation:
$\tfrac{(n-1)(n-2)}{2} = O(n^2)$ candidate moves.
Table~\ref{tab:scalability} projects runtime from the measured value at
$n = 200$ using the $O(n^2)$ scaling law.

\begin{table}[H]
\centering
\rowcolors{2}{gray!8}{white}
\begin{tabularx}{\linewidth}{cccX}
\toprule
$n$ & Time per iter.\ (proj.) & Total time (proj.) & Feasibility \\
\midrule
200     & baseline            & $<1$\,s        & Yes --- upper benchmark limit \\
500     & ${\sim}6\times$     & ${\sim}$few s  & Yes --- manageable \\
1\,000  & ${\sim}25\times$    & ${\sim}$minutes & Borderline \\
5\,000  & ${\sim}625\times$   & hours          & No --- too slow \\
10\,000 & ${\sim}2\,500\times$ & days          & No --- impractical \\
\bottomrule
\end{tabularx}
\caption{Projected scalability of Tabu Search (2-opt, \texttt{max\_iter}$=5\,000$)
  beyond the benchmark range.
  Time projections follow $T(n) = T(200) \times (n/200)^2$.}
\label{tab:scalability}
\end{table}

\paragraph{Conclusion.}
Tabu Search with a full 2-opt neighbourhood is practical for $n \leq 500$
and competitive for $n \leq 200$.
Beyond this range the $O(n^2)$ bottleneck dominates:

\begin{itemize}[leftmargin=1.5em]
  \item \textbf{$n \approx 500$--$1\,000$}: feasible but slow; a
    \textbf{candidate list} (restricting moves to the $k \approx 10$--$20$
    nearest neighbours) reduces per-iteration cost to $O(k \cdot n)$
    at a small quality trade-off.
  \item \textbf{$n \geq 5\,000$}: plain 2-opt TS is impractical.
    Lin--Kernighan--Helsgott (LKH-3) or population-based approaches
    (e.g.\ EAX) are necessary.
\end{itemize}

% =============================================================================
\section{Experimental Analysis}
\label{sec:experimental}
% =============================================================================

\paragraph{Scalability.}
Figure~\ref{fig:benchmark_summary} (centre) empirically confirms the $O(n^2)$
scaling predicted by complexity analysis.
Runtime grows approximately quadratically across the benchmark sizes.

\paragraph{Solution quality.}
Figure~\ref{fig:benchmark_summary} (left) shows that TS consistently improves
over the nearest-neighbour baseline across all sizes.
The gap is most pronounced for moderate $n$ (50--100), where 2-opt has
the most room to remove crossings.
For larger $n$ the greedy initial tour has more structure, and 5\,000
iterations may not be sufficient to fully exploit the neighbourhood.

\paragraph{Convergence behaviour.}
Figure~\ref{fig:benchmark_summary} (right) shows normalised convergence curves.
Early iterations yield steep improvements as the nearest-neighbour tour is
quickly refined; later iterations show diminishing returns.
The patience counter fires well before \texttt{max\_iter} for larger instances,
indicating that most gains are achieved within the first few hundred iterations.

\paragraph{Effect of tabu tenure $\theta$.}
The tabu tenure is the most sensitive hyperparameter:

\begin{center}
\rowcolors{2}{gray!8}{white}
\begin{tabularx}{\linewidth}{lX}
\toprule
\textbf{Regime} & \textbf{Observed behaviour} \\
\midrule
$\theta$ too small ($1$--$2$ for $n=100$) &
  Cycling: algorithm oscillates among a small set of solutions; converges
  quickly to a poor local optimum \\
$\theta$ too large ($> n/2$) &
  Over-restriction: admissible neighbourhood shrinks; poor moves are forced
  and the algorithm loses ability to exploit good regions \\
$\theta$ in sweet spot ($\lfloor\sqrt{n}\rfloor$--$n/5$) &
  Best balance between anti-cycling memory and search freedom;
  for $n=100$, values in $\theta \in [7, 20]$ are effective \\
\bottomrule
\end{tabularx}
\end{center}

A \textbf{dynamic tenure} that oscillates between a minimum and maximum value
prevents both pathological behaviours simultaneously.

\paragraph{Silent failures.}
Like the GNN solver, TS always outputs a tour with no internal quality
estimate.
A simple post-hoc check (compare tour length to the NN baseline) can flag
runs where TS failed to improve, enabling detection of poor solutions before
deployment.

\paragraph{Improvement directions.}

\begin{center}
\rowcolors{2}{gray!8}{white}
\begin{tabularx}{\linewidth}{lX}
\toprule
\textbf{Direction} & \textbf{Expected benefit} \\
\midrule
Candidate lists ($k$-NN) & Reduce per-iteration cost from $O(n^2)$ to
  $O(k n)$; enables $n > 1\,000$ \\
3-opt / Lin--Kernighan moves & Larger neighbourhood; better quality
  at modest extra cost \\
Random restarts & Escape large basins of attraction;
  best solution taken over multiple runs \\
Dynamic tenure & Avoids cycling and over-restriction simultaneously \\
TSPTW node features & Extend to time-windowed instances by adding
  time-window constraints to the move-acceptance rule \\
\bottomrule
\end{tabularx}
\end{center}

% =============================================================================
\section*{References}
% =============================================================================

\begin{itemize}[leftmargin=1.5em]
  \item Glover F. (1986).
    \textit{Future paths for integer programming and links to artificial
    intelligence}.
    Computers \& Operations Research, 13(5), 533--549.
  \item Glover F. (1989).
    \textit{Tabu search --- Part I}.
    ORSA Journal on Computing, 1(3), 190--206.
  \item Glover F. (1990).
    \textit{Tabu search --- Part II}.
    ORSA Journal on Computing, 2(1), 4--32.
  \item Glover F., Laguna M. (1997).
    \textit{Tabu Search}.
    Kluwer Academic Publishers.
  \item Gendreau M., Potvin J.-Y. (Eds.) (2010).
    \textit{Handbook of Metaheuristics} (2nd ed.), Chapter 2: Tabu Search.
    Springer.
  \item Lin S., Kernighan B.W. (1973).
    \textit{An effective heuristic algorithm for the traveling-salesman
    problem}.
    Operations Research, 21(2), 498--516.
  \item Vinyals O., Fortunato M., Jaitly N. (2015).
    \textit{Pointer Networks}.
    NeurIPS 2015.
  \item Kool W., Van Hoof H., Welling M. (2019).
    \textit{Attention, Learn to Solve Routing Problems!}
    ICLR 2019.
  \item Joshi C.K., Laurent T., Bresson X. (2019).
    \textit{An Efficient Graph Convolutional Network Technique for the
    Travelling Salesman Problem}.
    arXiv:1906.01227.
\end{itemize}

\end{document}
